\section{Open problems}\label{sec:open}

\paragraph{The full bias boundary.}
Does the bimodal interval keep widening, namely is $b_a$ strictly
decreasing on $(a_*,1)$? It does so near $a_*$
(Proposition~\ref{prop:bias-cusp}), and Figure~\ref{fig:phase}
suggests the same elsewhere. Theorem~\ref{thm:bias-continuity} gives
continuous differentiability and a unique minimizer for every $a>a_*$, while
Corollary~\ref{cor:bias-boundary-rate} gives
$b_a=O((1-a)^2)$ as $a\to1$. Is this the exact order, and is there
an asymptotic constant? The three-point bounds following
Theorem~\ref{thm:bias-continuity} apply uniformly when their hypotheses
are verified on parameter intervals. The missing step is a positive
parameter derivative at the maximizing points, with control near one.
More computed points do not settle what happens between them.

\paragraph{A finite-length threshold.}
For even $N$, does the central curvature have exactly one zero
$a_N^{\mathrm{crit}}$ in $(0,1)$ at every length?
Proposition~\ref{prop:finite-threshold-location} proves this for all
sufficiently large even $N$, on the whole memory interval.
At every even length a zero exists and all zeros lie below $a_*$.
The eventual unique zero increases to $a_*$ with leading shift of
order $N^{-(2a_*-1)}$.
The computed zero is near $0.6563$
for $N=4000$ and $0.6592$ for $N=8000$
(Remark~\ref{rem:modes-threshold}), still below $a_*$ in both cases.
The exponent $2a_*-1\approx0.3376$ explains the slow convergence.
The remaining issue is an effective cutoff and uniqueness at the
smaller lengths.

\paragraph{Between rare reversals and fixed memory.}
Corollary~\ref{cor:wider-mode-range} applies when
$x\to\infty$ and $x^3\log^2N/N\to0$.
The fixed-memory density tends to the same Landau profile as $a\to1$
(Theorem~\ref{thm:near-one-profile}), now with a rate.
Can we cover the scales between them with one approximation?
Our bounds leave room for improvement, judging by the computed errors
in Remark~\ref{rem:bulk}. The spatial derivative estimates in
\eqref{eq:profile-derivative-rate} describe the main conditional peak.
To find an asymptotic equivalent for $b_a$, we also need the much
smaller reflected contribution at the point where the other peak disappears.

\paragraph{The Drezner--Farnum model and higher dimensions.}
Drezner and Farnum's abstract says their family includes bimodal laws.
Changing its innovation probability $p$ changes the transition rule,
while changing the first-step bias $q$ mixes two reflected laws.
Our mixture argument therefore does not settle the shapes for $p\ne1/2$.
In higher dimensions, the fixed-point theory of~\cite{glr2026} gives
densities, but our sign argument needs points ordered along a line.
Even the meaning of a central depression needs care
(the usual axis-step walk is not rotationally symmetric).
We have not computed either extension.
